% ==============================================================================
% Figura 5: Máquina de Estados Finitos (FSM) del Firmware furuta_control.c (1 kHz)
% Diseño por construcción (Layout by Construction): posicionamiento relativo,
% anclas exactas, contenedores elásticos (fit) y cero cruces de transiciones.
% Bibliotecas TikZ requeridas: arrows.meta, positioning, fit, calc.
% ==============================================================================
\begin{figure*}[!t]
\centering
\resizebox{\textwidth}{!}{%
\begin{tikzpicture}[
    font=\sffamily,
    % --- Paleta semántica suave (ámbito local al entorno tikzpicture) ---
    /utils/exec={%
        \definecolor{fsmIdleFill}{HTML}{F1F5F9}%
        \definecolor{fsmIdleBdr}{HTML}{475569}%
        \definecolor{fsmIdleHdr}{HTML}{334155}%
        \definecolor{fsmServFill}{HTML}{EFF6FF}%
        \definecolor{fsmServBdr}{HTML}{2563EB}%
        \definecolor{fsmServHdr}{HTML}{1D4ED8}%
        \definecolor{fsmServZoneBdr}{HTML}{60A5FA}%
        \definecolor{fsmServZonePill}{HTML}{DBEAFE}%
        \definecolor{fsmServZoneTxt}{HTML}{1E3A8A}%
        \definecolor{fsmWarmFill}{HTML}{FFF7ED}%
        \definecolor{fsmWarmBdr}{HTML}{EA580C}%
        \definecolor{fsmWarmHdr}{HTML}{C2410C}%
        \definecolor{fsmWarmSubFill}{HTML}{FFEDD5}%
        \definecolor{fsmWarmSubBdr}{HTML}{FB923C}%
        \definecolor{fsmHeroFill}{HTML}{ECFDF5}%
        \definecolor{fsmHeroBdr}{HTML}{059669}%
        \definecolor{fsmHeroHdr}{HTML}{047857}%
        \definecolor{fsmAutoZoneBdr}{HTML}{0D9488}%
        \definecolor{fsmAutoZonePill}{HTML}{CCFBF1}%
        \definecolor{fsmAutoZoneTxt}{HTML}{115E59}%
        \definecolor{fsmStopFill}{HTML}{FEF2F2}%
        \definecolor{fsmStopBdr}{HTML}{DC2626}%
        \definecolor{fsmStopHdr}{HTML}{B91C1C}%
        \definecolor{fsmAlertFill}{HTML}{FFF1F2}%
        \definecolor{fsmAlertBdr}{HTML}{E11D48}%
        \definecolor{fsmAlertTxt}{HTML}{9F1239}%
        \definecolor{fsmInk}{HTML}{1E293B}%
        \definecolor{fsmMuted}{HTML}{475569}%
    },
    % --- Estilos de nodos de estado (Layout by Construction) ---
    fsmNode/.style={
        rectangle,
        rounded corners=8pt,
        align=center,
        inner xsep=7pt,
        inner ysep=6pt,
        text=fsmInk,
        line width=1.0pt
    },
    stateIdle/.style={
        fsmNode,
        fill=fsmIdleFill,
        draw=fsmIdleBdr,
        line width=1.25pt,
        text width=3.95cm,
        minimum height=2.40cm
    },
    stateServ/.style={
        fsmNode,
        fill=fsmServFill,
        draw=fsmServBdr,
        text width=4.40cm,
        minimum height=1.65cm
    },
    stateSwing/.style={
        fsmNode,
        fill=fsmWarmFill,
        draw=fsmWarmBdr,
        line width=1.1pt,
        text width=4.75cm,
        minimum height=2.35cm
    },
    stateHero/.style={
        fsmNode,
        fill=fsmHeroFill,
        draw=fsmHeroBdr,
        line width=1.75pt,
        text width=5.25cm,
        minimum height=2.35cm
    },
    stateRecover/.style={
        fsmNode,
        fill=fsmWarmFill,
        draw=fsmWarmBdr,
        line width=1.1pt,
        text width=11.4cm,
        minimum height=2.05cm
    },
    stateStop/.style={
        fsmNode,
        fill=fsmStopFill,
        draw=fsmStopBdr,
        line width=1.1pt,
        text width=8.2cm,
        minimum height=1.40cm
    },
    % --- Estilos de transiciones (flechas canónicas) ---
    transMain/.style={
        -{Stealth[length=2.7mm, width=1.95mm]},
        line width=1.15pt,
        draw=fsmAutoZoneTxt,
        shorten >=1.5pt,
        shorten <=1.0pt
    },
    transWarm/.style={
        -{Stealth[length=2.7mm, width=1.95mm]},
        line width=1.15pt,
        draw=fsmWarmHdr,
        shorten >=1.5pt,
        shorten <=1.0pt
    },
    transServ/.style={
        -{Stealth[length=2.3mm, width=1.7mm]},
        line width=0.95pt,
        draw=fsmServHdr,
        shorten >=1.5pt,
        shorten <=1.0pt
    },
    transRet/.style={
        -{Stealth[length=2.3mm, width=1.7mm]},
        line width=0.90pt,
        draw=fsmIdleBdr,
        shorten >=1.5pt,
        shorten <=1.0pt
    },
    transBypass/.style={
        -{Stealth[length=2.5mm, width=1.8mm]},
        line width=1.05pt,
        draw=fsmHeroHdr,
        dashed,
        dash pattern=on 5pt off 3pt,
        shorten >=1.5pt,
        shorten <=1.0pt
    },
    transStop/.style={
        -{Stealth[length=2.6mm, width=1.9mm]},
        line width=1.1pt,
        draw=fsmStopHdr,
        shorten >=1.5pt,
        shorten <=1.0pt
    },
    % --- Etiquetas de transición ---
    elbl/.style={
        font=\small\sffamily,
        text=fsmInk,
        align=center,
        inner sep=1.8pt
    },
    elblBox/.style={
        elbl,
        fill=white,
        rounded corners=2pt,
        inner xsep=3.5pt,
        inner ysep=2pt
    },
    % --- Badges y contenedores ---
    zoneServ/.style={
        draw=fsmServZoneBdr,
        dashed,
        line width=0.95pt,
        rounded corners=10pt,
        fill=none,
        inner xsep=0.34cm,
        inner ysep=0.45cm
    },
    zoneAuto/.style={
        draw=fsmAutoZoneBdr,
        dashed,
        line width=1.1pt,
        rounded corners=10pt,
        fill=none,
        inner xsep=0.36cm,
        inner ysep=0.50cm
    },
    zonePill/.style={
        rectangle,
        rounded corners=3.5pt,
        font=\small\bfseries\sffamily,
        inner xsep=6pt,
        inner ysep=2.5pt
    }
]

% ==============================================================================
% 1. COLOCACIÓN RELATIVA DE ESTADOS (LAYOUT BY CONSTRUCTION)
% ==============================================================================

% --- Núcleo central: Estado Inicial IDLE (Estado 0) ---
\node[stateIdle] (idle) at (0,0) {%
    \colorbox{fsmIdleHdr}{\textcolor{white}{\footnotesize\bfseries Est.\,0}}\\[3pt]
    {\large\bfseries IDLE}\\[2pt]
    {\small Estado inicial / reposo}\\
    {\small Rueda libre (\texttt{L298N\_Coast})}\\
    {\small \mbox{$u = \qty{0}{\volt}$}}
};

% --- Rama Izquierda: Servicio y Diagnóstico (Estados 1, 2, 3–4) ---
\node[stateServ, left=3.75cm of idle] (char) {%
    \colorbox{fsmServHdr}{\textcolor{white}{\footnotesize\bfseries Est.\,2}}\\[3pt]
    {\normalsize\bfseries CHARACTERIZATION}\\[2pt]
    {\small Ensayos \texttt{CHAR1}--\texttt{CHAR4}}\\
    {\small Validaci\'on PID (\texttt{CHAR5A/B})}
};

\node[stateServ, above=0.95cm of char] (cal) {%
    \colorbox{fsmServHdr}{\textcolor{white}{\footnotesize\bfseries Est.\,1}}\\[3pt]
    {\normalsize\bfseries CALIBRATING}\\[2pt]
    {\small Promedio de 64 muestras}\\
    {\small Cero de encoders AS5600}
};

\node[stateServ, below=0.95cm of char] (pidmp) {%
    \colorbox{fsmServHdr}{\textcolor{white}{\footnotesize\bfseries Est.\,3--4}}\\[3pt]
    {\small\bfseries MOTOR\_PID\_VEL / POS}\\[2pt]
    {\small PID Tustin (vel.) / Servo}\\
    {\small Clegg + freno activo}
};

% --- Rama Principal Derecha: Ciclo Autónomo de Control Furuta ---
% Fila superior del ciclo: SWING_UP (Est. 5) y BALANCE_PID (Est. 6, HERO)
\node[stateSwing, right=3.80cm of idle, yshift=1.70cm] (swing) {%
    \colorbox{fsmWarmHdr}{\textcolor{white}{\footnotesize\bfseries Est.\,5}}\\[3pt]
    {\large\bfseries SWING\_UP}\\[2pt]
    {\small Bombeo resonante\\en cuadratura}\\
    {\small Carrera adaptativa: $50^\circ / 36^\circ / 22^\circ$}
};

\node[stateHero, right=2.90cm of swing] (balance) {%
    \colorbox{fsmHeroHdr}{\textcolor{white}{\footnotesize\bfseries Est.\,6 \;$\bullet$\; EQUILIBRIO}}\\[3pt]
    {\large\bfseries\color{fsmHeroHdr} BALANCE\_PID}\\[2pt]
    {\small M4 Ultra-Silk + Dual-Notch\\(\qty{11}{\hertz} / \qty{6.2}{\hertz})}\\
    {\small Slew adaptativo; sin límite angular $\theta$}
};

% Fila inferior del ciclo: RECOVER_HOME (Est. 8) centrado bajo SWING_UP y BALANCE_PID
\coordinate (auto_row1_mid) at ($(swing.south west)!0.5!(balance.south east)$);
\node[stateRecover, anchor=north] (recover) at ($(auto_row1_mid)+(0,-1.95cm)$) {%
    \colorbox{fsmWarmHdr}{\textcolor{white}{\footnotesize\bfseries Est.\,8}}\quad
    {\large\bfseries RECOVER\_HOME}\quad
    {\small (Auto-recuperación en 3 subfases)}\\[5pt]
    \fcolorbox{fsmWarmSubBdr}{fsmWarmSubFill}{\small\strut\textbf{F0:} Rueda libre (\qty{300}{\milli\second})}%
    \;\mbox{$\rightarrow$}\;%
    \fcolorbox{fsmWarmSubBdr}{fsmWarmSubFill}{\small\strut\textbf{F1:} Retorno suave a $\theta = 0^\circ$}\\[4pt]%
    \mbox{$\rightarrow$}\;%
    \fcolorbox{fsmWarmSubBdr}{fsmWarmSubFill}{\small\strut\textbf{F2:} Freno activo hasta calma (\qty{250}{\milli\second})}
};

% ==============================================================================
% 2. CONTENEDORES ELÁSTICOS (FIT) CON ENCABEZADOS TIPO PILL
% ==============================================================================

% Contenedor Zona Servicio y Diagnóstico
\coordinate (serv_top_pad) at ($(cal.north)+(0,0.20cm)$);
\node[zoneServ, fit=(cal)(char)(pidmp)(serv_top_pad)] (serv_zone) {};
\node[zonePill, fill=fsmServZonePill, draw=fsmServZoneBdr, text=fsmServZoneTxt]
    at (serv_zone.north) {Servicio y Diagnóstico};

% Contenedor Superestado Ciclo Autónomo Cerrado
\coordinate (auto_top_pad) at ($(swing.north)+(0,0.30cm)$);
\node[zoneAuto, fit=(swing)(balance)(recover)(auto_top_pad)] (auto_zone) {};
\node[zonePill, fill=fsmAutoZonePill, draw=fsmAutoZoneBdr, text=fsmAutoZoneTxt, anchor=west]
    at ([xshift=0.35cm]auto_zone.north west) {Ciclo Autónomo Cerrado de Control Furuta};

% --- Parada Ordenada: STOP_HOME (Est. 9) bajo el superestado autónomo ---
\node[stateStop, below=1.10cm of auto_zone.south] (stophome) {%
    \colorbox{fsmStopHdr}{\textcolor{white}{\footnotesize\bfseries Est.\,9}}\quad
    {\large\bfseries STOP\_HOME}\\[3pt]
    {\small Parada ordenada (botón \textbf{Detener} MATLAB):\\retorno suave a \mbox{$\theta = 0^\circ$} y paso a \texttt{IDLE}}
};

% ==============================================================================
% 3. TRANSICIONES Y ETIQUETAS (CERO CRUCES, ANCLAS EXACTAS)
% ==============================================================================

% --- Entrada inicial a IDLE (Inicio / ISR 1 kHz) ---
\coordinate (init_dot) at ([xshift=-0.70cm, yshift=0.95cm]idle.north);
\fill[fsmIdleHdr] (init_dot) circle (3.2pt);
\draw[transRet, line width=1.15pt, draw=fsmIdleHdr]
    (init_dot) -- ([xshift=-0.70cm]idle.north);
\node[elbl, above=2.5pt] at (init_dot) {\textbf{Reset / Inicio}\\ISR $\qty{1}{\kilo\hertz}$};

% --- Carriles de ruteo ortogonal en la brecha izquierda (Service <-> IDLE) ---
\coordinate (laneL_near_idle) at ($(idle.west)+(-0.38cm,0)$);
\coordinate (laneL_mid)       at ($(idle.west)+(-0.76cm,0)$);
\coordinate (gapL_mid_near)   at ($(serv_zone.east)!0.5!(laneL_near_idle)$);
\coordinate (gapL_mid_lane)   at ($(serv_zone.east)!0.5!(laneL_mid)$);

% IDLE -> CALIBRATING (ZERO / Botón B1)
\coordinate (cal_in)  at ([yshift=0.22cm]cal.east);
\coordinate (cal_out) at ([yshift=-0.22cm]cal.east);
\draw[transServ, rounded corners=6pt]
    ([yshift=0.88cm]idle.west) -- ([yshift=0.88cm]idle.west -| laneL_near_idle)
    -- (cal_in -| laneL_near_idle)
    -- (cal_in);
\node[elbl, above=2pt] at (cal_in -| gapL_mid_near) {\texttt{ZERO} / Botón B1};

% CALIBRATING -> IDLE (Calibración completa)
\draw[transRet, rounded corners=6pt]
    (cal_out) -- (cal_out -| laneL_mid)
    -- ([yshift=0.58cm]idle.west -| laneL_mid)
    -- ([yshift=0.58cm]idle.west);
\node[elbl, below=2pt, text=fsmMuted] at (cal_out -| gapL_mid_lane) {Calibrado};

% IDLE <-> CHARACTERIZATION (paralelas rectas)
\draw[transServ]
    ([yshift=0.20cm]idle.west) -- ([yshift=0.20cm]char.east);
\node[elbl, above=2pt] at ([yshift=0.20cm]gapL_mid_lane) {\texttt{CHAR1}--\texttt{CHAR5}};
\draw[transRet]
    ([yshift=-0.20cm]char.east) -- ([yshift=-0.20cm]idle.west);
\node[elbl, below=2pt, text=fsmMuted] at ([yshift=-0.20cm]gapL_mid_lane) {\texttt{test\_finished}};

% IDLE -> MOTOR_PID_VEL / POS
\draw[transServ, rounded corners=6pt]
    ([yshift=-0.82cm]idle.west) -- ([yshift=-0.82cm]idle.west -| laneL_near_idle)
    -- (pidmp.east -| laneL_near_idle)
    -- (pidmp.east);
\node[elbl, above=2pt] at (pidmp.east -| gapL_mid_near) {\texttt{VEL=} / \texttt{POS=}};

% --- Carril de ruteo derecho (IDLE -> Ciclo Autónomo) ---
\coordinate (laneR_near_idle) at ($(idle.east)+(0.42cm,0)$);
\coordinate (gapR_mid)        at ($(laneR_near_idle)!0.5!(auto_zone.west)$);

% IDLE -> SWING_UP (START/SWING si está en 0° y reposo)
\draw[transMain, rounded corners=6pt]
    ([yshift=0.65cm]idle.east) -- ([yshift=0.65cm]idle.east -| laneR_near_idle)
    -- (swing.west -| laneR_near_idle)
    -- (swing.west);
\node[elbl, above=2pt] at (swing.west -| gapR_mid) {\texttt{START} / \texttt{SWING}};
\node[elbl, below=2pt, text=fsmMuted] at (swing.west -| gapR_mid) {(en $0^\circ$ y reposo)};

% IDLE -> RECOVER_HOME (START/SWING si brazo != 0° o en movimiento)
\draw[transWarm, rounded corners=6pt]
    ([yshift=-0.65cm]idle.east) -- ([yshift=-0.65cm]idle.east -| laneR_near_idle)
    -- (recover.west -| laneR_near_idle)
    -- (recover.west);
\node[elbl, above=2pt] at (recover.west -| gapR_mid) {\texttt{START} / \texttt{SWING}};
\node[elbl, below=2pt, text=fsmMuted] at (recover.west -| gapR_mid) {($\theta \neq 0^\circ$ o mov.)};

% IDLE -> BALANCE_PID (acceso directo por corredor superior, cero cruces)
\coordinate (top_rail) at ($(auto_zone.north)+(0,0.75cm)$);
\coordinate (bal_exit) at ([xshift=0.70cm]idle.north);
\draw[transBypass, rounded corners=7pt]
    (bal_exit) -- (bal_exit |- top_rail)
    -- node[pos=0.55, above=2pt, elbl, text=fsmHeroHdr]
       {\texttt{BALANCE} \;{\color{fsmMuted}(acceso directo, \texttt{auto\_swingup=false})}}
       (balance.north |- top_rail)
    -- (balance.north);

% --- Transiciones internas del Ciclo Autónomo Cerrado ---

% SWING_UP -> BALANCE_PID (Transferencia sin saltos / bumpless)
\draw[transMain, draw=fsmHeroHdr, line width=1.35pt]
    (swing.east) --
    node[above=3pt, elbl] {$|\tilde{\alpha}| < 20^\circ$\\$|\dot{\alpha}| < \qty{4.8}{\radian\per\second}$}
    node[below=3pt, elbl, text=fsmHeroHdr] {\bfseries Transferencia\\[-1pt]\bfseries sin saltos}
    (balance.west);

% SWING_UP -> RECOVER_HOME (Timeout > 5.5 s)
\coordinate (sw_tmo_top) at ([xshift=-1.15cm]swing.south);
\coordinate (sw_tmo_bot) at (sw_tmo_top |- recover.north);
\draw[transWarm]
    (sw_tmo_top) -- node[pos=0.50, elblBox, text=fsmWarmHdr] {Timeout\\$t > \qty{5.5}{\second}$} (sw_tmo_bot);

% RECOVER_HOME -> SWING_UP (Relanza automáticamente SWING_UP — ciclo cerrado)
\coordinate (sw_rel_top) at ([xshift=1.05cm]swing.south);
\coordinate (sw_rel_bot) at (sw_rel_top |- recover.north);
\draw[transMain, line width=1.25pt]
    (sw_rel_bot) -- node[pos=0.50, right=4pt, elbl, align=left, text=fsmAutoZoneTxt]
        {\textbf{Ciclo cerrado}\\Calma $\qty{250}{\milli\second}$\\$\rightarrow$ relanza \texttt{SWING\_UP}} (sw_rel_top);

% BALANCE_PID -> RECOVER_HOME (Watchdog de caída |alpha_tilde| > 32° y auto_swingup=true)
\coordinate (bal_wd_bot) at (balance.south |- recover.north);
\draw[transWarm, line width=1.25pt]
    (balance.south) -- node[pos=0.50, elblBox, text=fsmWarmHdr]
        {\textbf{Watchdog caída}\\$|\tilde{\alpha}| > 32^\circ \land \texttt{auto\_swingup=true}$} (bal_wd_bot);

% --- Parada ordenada: Ciclo Autónomo -> STOP_HOME -> IDLE ---
\draw[transStop]
    (auto_zone.south) -- node[midway, right=5pt, elbl, align=left, text=fsmStopHdr]
        {\texttt{STOPHOME} {\color{fsmMuted}(desde ciclo autónomo)}} (stophome.north);

\draw[transStop, rounded corners=7pt]
    (stophome.west) -| node[pos=0.25, above=2pt, elbl, text=fsmStopHdr]
        {Brazo en $0^\circ$ ($\qty{150}{\milli\second}$) $\rightarrow$ \texttt{L298N\_Coast}} (idle.south);

% ==============================================================================
% 4. NOTA GLOBAL DE EMERGENCIA (CUADRO DISCRETO DE ANCHO COMPLETO)
% ==============================================================================
\coordinate (banner_y) at ($(stophome.south)+(0,-0.55cm)$);
\coordinate (banner_L) at (serv_zone.west |- banner_y);
\coordinate (banner_R) at (auto_zone.east |- banner_y);
\coordinate (banner_mid) at ($(banner_L)!0.5!(banner_R)$);

\path let \p1 = (serv_zone.west), \p2 = (auto_zone.east),
          \n1 = {\x2 - \x1 - 16pt - 1.9pt} in
    node[
        rectangle,
        rounded corners=6pt,
        draw=fsmAlertBdr,
        fill=fsmAlertFill,
        line width=0.95pt,
        anchor=north,
        inner xsep=8pt,
        inner ysep=5pt,
        align=center,
        font=\small\sffamily,
        text=fsmInk,
        text width=\n1
    ] (global_note) at (banner_mid) {%
        \colorbox{fsmAlertBdr}{\textcolor{white}{\footnotesize\bfseries PARADA GLOBAL DE EMERGENCIA}}\\[3pt]
        El comando UART \texttt{STOP} o la pulsaci\'on del \textbf{Bot\'on B1} (\texttt{PC13}) fuerzan el retorno inmediato a \texttt{IDLE} (\mbox{$u=\qty{0}{\volt}$}) desde \textbf{cualquier estado}.\\[2pt]
        {\color{fsmMuted}(Si \mbox{$|\tilde{\alpha}|>32^\circ$} en \texttt{BALANCE\_PID} con \texttt{auto\_swingup=false}, tambi\'en conmuta autom\'aticamente a \texttt{IDLE}.)}
    };

\end{tikzpicture}%
}%
\caption{Máquina de estados finitos (FSM) del firmware de control digital (\texttt{furuta\_control.c}, lazo isócrono \texttt{TIM6} a \qty{1.0}{\kilo\hertz}). Se muestran el estado seguro \texttt{IDLE} (estado~0), la rama de servicio y diagnóstico (estados~1--4), el ciclo autónomo cerrado (\texttt{SWING\_UP}, \texttt{BALANCE\_PID} y \texttt{RECOVER\_HOME}, estados~5, 6 y 8), la parada ordenada \texttt{STOP\_HOME} (estado~9) y la parada global de emergencia (\texttt{STOP} o botón \texttt{B1}).}
\label{fig:fsm_firmware}
\end{figure*}
